\chapter{Half-Life, Decay and Stability}\label{ch:rs:half-life-decay-and-stability}

McLean and Pontiff (2016) followed 97 variables that published studies had shown to predict the cross-section of stock returns,
past the samples of the studies and past their publication. Portfolio returns were 26\% lower out of sample and 58\% lower after
publication; they read the difference, 32\%, as the effect of investors trading on what they had read. Every signal a desk owns is
on some such path, and the questions of this chapter are how fast a \omterm{def:rs:anatomy-of-a-predictor:predictor}{predictor}'s information runs out as the horizon lengthens, how
stable it is from year to year, and how to tell a \omterm{def:rs:anatomy-of-a-predictor:predictor}{predictor} that has died from one that was unlucky. The clock of the first question is
days; of the second, years; the third is a problem of statistical power, and its answer is often that the data cannot tell.

\section{Decay across horizons}

\begin{definition}[IC decay curve, signal autocorrelation, signal half-life]\label{def:rs:half-life-decay-and-stability:decay}
The \emph{IC decay curve}\index{IC decay curve} of a \omterm{def:rs:anatomy-of-a-predictor:predictor}{predictor} is its \omterm{def:rs:anatomy-of-a-predictor:ic}{information coefficient} with the return of the single period
$t + h$, as a function of the lag $h$. The \emph{signal autocorrelation}\index{signal autocorrelation} is the cross-sectional rank
correlation of the \omterm{def:rs:anatomy-of-a-predictor:predictor}{predictor} with itself $k$ periods earlier, averaged over dates. The \emph{signal half-life}\index{signal half-life} is
the lag at which it falls to one half; for a first-order autoregression with coefficient $\rho$ per period it is $-\ln 2/\ln\rho$ (the
half-life of Book~4, chapter~4).
\end{definition}

The two curves answer different questions. The \omterm{def:rs:half-life-decay-and-stability:decay}{IC decay curve} says how long the information lasts, and so the holding period; the \omterm{def:rs:half-life-decay-and-stability:decay}{signal
autocorrelation} says how fast the \omterm{def:rs:anatomy-of-a-predictor:predictor}{predictor} itself changes, and so the turnover. A rank-weighted book rebuilt every period trades a fraction
of about $1 - \rho$ of itself, where $\rho$ is the lag-one \omterm{def:rs:half-life-decay-and-stability:decay}{signal autocorrelation}. On \texttt{firm.synthmkt}
(\cref{fig:rs:half-life-decay-and-stability:curves} and the table), four \omterm{def:rs:anatomy-of-a-predictor:predictor}{predictors} show four shapes:

\begin{center}
\small
\begin{tabular}{@{}lccccc@{}}
\toprule
\omterm{def:rs:anatomy-of-a-predictor:predictor}{predictor} & IC, $h = 1$ & curve half-life & autocorr. & \omterm{def:rs:half-life-decay-and-stability:decay}{signal half-life} & turnover\\
\midrule
reversal, 1 day & 0.037 & $<$ 1 day & $-0.04$ & --- & 104\%\\
surprise, 60 days & 0.010 & 20 days (fit) & 0.983 & 39 days & 1.7\%\\
momentum 12-1 & 0.008 & none in 60 days & 0.995 & 130 days & 0.53\%\\
book-to-price & 0.004 & none in 60 days & 0.9995 & 1\,386 days & 0.05\%\\
\bottomrule
\end{tabular}
\end{center}

\begin{omfigure}
\begin{tikzpicture}
\begin{axis}[width=11cm,height=5.6cm,xlabel={lag $h$ of the single-day return (days)},ylabel={mean rank IC},
  tick label style={font=\small},label style={font=\small},xmin=0,xmax=61,ymin=-0.01,ymax=0.04,grid=major,grid style={gray!25},
  yticklabel style={/pgf/number format/fixed,/pgf/number format/precision=2},scaled y ticks=false,legend columns=2,
  legend style={font=\footnotesize,at={(0.5,-0.25)},anchor=north,draw=none,fill=none,
  /tikz/every even column/.append style={column sep=8pt}},table/col sep=comma]
\addplot[omThm,thick,mark=*,mark size=1.4pt] table[x=h,y=reversal]{figdata/research/13-half-life-decay-and-stability/decay.csv};
\addplot[omDef,very thick,mark=square*,mark size=1.4pt] table[x=h,y=surprise]{figdata/research/13-half-life-decay-and-stability/decay.csv};
\addplot[black!60,thick,dashed,mark=triangle*,mark size=1.4pt] table[x=h,y=momentum]{figdata/research/13-half-life-decay-and-stability/decay.csv};
\addplot[gray,thick,dotted,mark=o,mark size=1.4pt] table[x=h,y=bp]{figdata/research/13-half-life-decay-and-stability/decay.csv};
\legend{one-day reversal,earnings surprise (held 60 days),momentum 12-1,book-to-price}
\end{axis}
\end{tikzpicture}

\omcaption{\omterm{def:rs:half-life-decay-and-stability:decay}{IC decay curves} on \texttt{firm.synthmkt}: the rank IC of each \omterm{def:rs:anatomy-of-a-predictor:predictor}{predictor} known at the close of day $t$ with the return of day
$t + h$ alone, averaged over eight years of dates. Reversal lives one day; the surprise's planted 60-day drift falls to zero at 60 days;
momentum and value are too slow to decay within the window.}\label{fig:rs:half-life-decay-and-stability:curves}
\end{omfigure}

The table carries two warnings. The surprise's curve is not exponential: every event drifts at a constant rate for 60 days, and a signal
observed at a random age has a remaining drift that falls linearly to zero, so the curve is a straight line with a kink; an exponential
fit reports a half-life of 20 days for a planted life of 60. Fitting a form assumes it. And the slow \omterm{def:rs:anatomy-of-a-predictor:predictor}{predictors}' daily ICs, 0.004 to 0.008,
are close to their sampling noise; their decay cannot be measured on single days, only on longer returns and many years. The practical
reading is on the turnover column: the one-day reversal trades its whole book every day and pays costs to match (chapter~27), book-to-price
a twentieth of a percent.

\section{Stability across years}

\begin{definition}[IC stability]\label{def:rs:half-life-decay-and-stability:stability}
The \emph{IC stability}\index{IC stability} of a \omterm{def:rs:anatomy-of-a-predictor:predictor}{predictor} is the behaviour of its \omterm{def:rs:anatomy-of-a-predictor:ic}{information coefficient} across non-overlapping
periods: the dispersion of yearly (or monthly) ICs around their mean, compared with the dispersion their sampling error alone would produce.
\end{definition}

The monthly IC of a signal is a noisy number, and its noise is not only sampling error in the cross-section: when the \omterm{def:rs:anatomy-of-a-predictor:predictor}{predictor} loads on a
common factor, the factor's return that month moves every stock's return together and with it the IC. Momentum on the synthetic market
(planted constant for ten years) has a mean monthly IC of 0.023 and a standard deviation of 0.247; its yearly means run from $-0.11$ to
0.18, and a block-bootstrap 95\% interval for the mean runs from $-0.035$ to 0.071, across zero. The \omterm{def:rs:fundamental-analyst-and-event-features:sue}{earnings surprise} has a mean of 0.040
with a standard deviation of 0.041, and its yearly means stay between 0.02 and 0.06. Two constant \omterm{def:rs:anatomy-of-a-predictor:predictor}{predictors}, one of which looks stable and
one of which looks like a sequence of regimes: the difference is the factor exposure, not the truth.

\section{Structural breaks}

\begin{definition}[Structural break, CUSUM test]\label{def:rs:half-life-decay-and-stability:break}
A \emph{structural break}\index{structural break} is a change, at some date, in the parameters of the process generating a series (here,
the mean of a \omterm{def:rs:anatomy-of-a-predictor:predictor}{predictor}'s IC). The \emph{CUSUM test}\index{CUSUM test} of Brown, Durbin and Evans (1975) cumulates the standardised
recursive residuals (each observation minus the mean of those before it, rescaled) and rejects constancy when the path leaves a pair of
lines that widen linearly with time, $\pm a(\sqrt K + 2k/\sqrt K)$ with $a = 0.948$ at 5\% for $K$ residuals.
\end{definition}

The alternative is to try every date: fit one mean before and one after, compute the F statistic of the split, and take the largest over
the middle of the sample (Andrews, 1993, gave the theory of this sup-F statistic for a change at an unknown point; Bai and Perron, 1998,
extended it to several breaks). \texttt{firm.decay} computes its p-value by permutation, which keeps the series' own distribution (heavy
tails included) and assumes only that, without a break, the order of the months is irrelevant.

A planted break makes the comparison. With \texttt{pead\_break} at the start of year 7, the synthetic market's post-earnings drift is
removed entirely: the surprise's mean monthly IC falls from 0.032 before to 0.005 after (\cref{fig:rs:half-life-decay-and-stability:break}).
The sup-F statistic, 12.5, peaks 53 months into the series (the break is at 48, and old announcements keep drifting for up to 60 days after
it), with a permutation p-value of 0.013. The CUSUM does not cross its boundary: its path reaches 0.69 of it. With the drift cut by 70\% instead of removed, neither test sees anything in this market (sup-F p-value 0.93).
One market is one draw; over ten markets with different seeds, sup-F detects the removal in all ten, the 70\% cut in seven and a change that
never happened in one, the CUSUM in five, three and one. And momentum, which nothing changed, fails both: its CUSUM crosses in
the last month, its sup-F p-value is 0.023, after a few strong years. A 5\% test rejects a stable \omterm{def:rs:anatomy-of-a-predictor:predictor}{predictor} one time in twenty by design,
more often when the IC's variance changes over time, and a desk that monitors fifty \omterm{def:rs:anatomy-of-a-predictor:predictor}{predictors} monthly will see breaks every year that are
not there.

\begin{omfigure}
\begin{tikzpicture}
\begin{axis}[width=11cm,height=5.4cm,xlabel={year of the synthetic market},ylabel={12-month mean IC},tick label style={font=\small},
  label style={font=\small},xmin=2.8,xmax=10,ymin=-0.02,ymax=0.08,grid=major,grid style={gray!25},
  yticklabel style={/pgf/number format/fixed,/pgf/number format/precision=2},scaled y ticks=false,legend columns=3,
  legend style={font=\footnotesize,at={(0.5,-0.25)},anchor=north,draw=none,fill=none,
  /tikz/every even column/.append style={column sep=8pt}},
  table/col sep=comma]
\addplot[black!55,thick,dashed] table[x=year,y=intact]{figdata/research/13-half-life-decay-and-stability/rolling_ic.csv};
\addplot[omDef,very thick] table[x=year,y=broken]{figdata/research/13-half-life-decay-and-stability/rolling_ic.csv};
\draw[omThm,thick] (axis cs:6,-0.02) -- (axis cs:6,0.08) node[pos=0.92,left,font=\footnotesize]{planted break};
\legend{drift intact,drift removed from day 1\,512}
\end{axis}
\end{tikzpicture}

\omcaption{Twelve-month rolling mean of the \omterm{def:rs:fundamental-analyst-and-event-features:sue}{earnings surprise}'s monthly IC on \texttt{firm.synthmkt}, with the post-earnings drift intact and
with it removed from the start of year 7 (day 1\,512). The sup-F test finds the break (p-value 0.013); the CUSUM does not; in this market a 70\% cut is found by neither.}\label{fig:rs:half-life-decay-and-stability:break}
\end{omfigure}

\section{Post-publication decay}

\begin{definition}[Post-publication decay]\label{def:rs:half-life-decay-and-stability:postpub}
The \emph{post-publication decay}\index{post-publication decay} of a \omterm{def:rs:anatomy-of-a-predictor:predictor}{predictor} is the fall of its return, or its IC, after the study that
documented it was published, relative to the period the study could have examined.
\end{definition}

Three of the most studied factors show it in the Kenneth French data library (the chapter stores only derived statistics).
Taking each from July 1963 to the end of the year its paper appeared (Banz, 1981, for size; Fama and French, 1992, for value; Jegadeesh and
Titman, 1993, for momentum) and from then to July 2026:

\begin{center}
\small
\begin{tabular}{@{}lccccc@{}}
\toprule
factor (paper) & mean a month, before & $t$ & mean a month, after & $t$ & decline\\
\midrule
SMB (1981) & 0.47\% & 2.25 & 0.01\% & 0.10 & 97\%\\
HML (1992) & 0.42\% & 3.12 & 0.19\% & 1.16 & 55\%\\
Mom (1993) & 0.85\% & 4.76 & 0.37\% & 1.54 & 56\%\\
\bottomrule
\end{tabular}
\end{center}

The declines bracket McLean and Pontiff's average, and the rolling ten-year means show them
(\cref{fig:rs:half-life-decay-and-stability:factors}). Yet a sup-F test for one break over the whole monthly history from 1963 rejects for
none of the three at 5\% (permutation p-values 0.18, 0.15 and 0.053). Factor returns are noisy enough that a halving over decades is hard
to date statistically, and a publication date is only one of many candidate causes: Chordia, Subrahmanyam and Tong (2014) found that the
returns of a portfolio of prominent anomalies roughly halved after decimalisation, with hedge-fund assets, short interest and turnover
behind the decline. A replication adds a third reading: Hou, Xue and Zhang (2020) found that 65\% of 452 anomalies could not clear a $t$
statistic of 1.96 once microcaps were mitigated, and that even the replicated ones were much smaller than originally reported. Some of the
decay is the original sample's selection.

\begin{omfigure}
\begin{tikzpicture}
\begin{axis}[width=11cm,height=5.6cm,xlabel={year},ylabel={120-month mean (\% a month)},tick label style={font=\small},
  label style={font=\small},xmin=1972,xmax=2027,ymin=-0.5,ymax=1.5,grid=major,grid style={gray!25},
  xticklabel style={/pgf/number format/1000 sep=},legend columns=3,
  legend style={font=\footnotesize,at={(0.5,-0.25)},anchor=north,draw=none,fill=none,
  /tikz/every even column/.append style={column sep=8pt}},table/col sep=comma]
\addplot[gray,thick] table[x=year,y=SMB]{figdata/research/13-half-life-decay-and-stability/factors.csv};
\addplot[omDef,thick] table[x=year,y=HML]{figdata/research/13-half-life-decay-and-stability/factors.csv};
\addplot[omThm,very thick] table[x=year,y=Mom]{figdata/research/13-half-life-decay-and-stability/factors.csv};
\draw[gray,dashed] (axis cs:1981,-0.5) -- (axis cs:1981,1.5);
\draw[omDef,dashed] (axis cs:1992,-0.5) -- (axis cs:1992,1.5);
\draw[omThm,dashed] (axis cs:1993,-0.5) -- (axis cs:1993,1.5);
\legend{SMB (published 1981),HML (1992),Mom (1993)}
\end{axis}
\end{tikzpicture}

\omcaption{Rolling 120-month mean monthly return of the US size, value and momentum factors at each December from 1973, with the years of the
papers that published them (dashed). Derived from the Kenneth R. French Data Library (monthly factors, 202607 CRSP file).}\label{fig:rs:half-life-decay-and-stability:factors}
\end{omfigure}

\section{Measuring the half-life of a predictor card}

Every \omterm{def:rs:anatomy-of-a-predictor:card}{predictor card} (chapter~6) carries a horizon and half-life field, and this chapter's tools fill it: the \omterm{def:rs:half-life-decay-and-stability:decay}{IC decay curve} at the card's
sampling frequency, its fitted half-life with a bootstrap interval (and the warning of the surprise: plot the curve before trusting the
form), the \omterm{def:rs:half-life-decay-and-stability:decay}{signal autocorrelation} and the turnover it implies, and the monthly IC series with its CUSUM and sup-F tests, refreshed every
month. The hardest question is the last one: when the recent IC is half the historical one, is the \omterm{def:rs:anatomy-of-a-predictor:predictor}{predictor} dying? Two periods of $n$ months
each with a monthly IC of mean $\mu$ and standard deviation $\sigma$ give two means with standard error $\sigma/\sqrt n$; the probability of
measuring a halving when nothing changed and when the IC truly halved depends only on the $t$ statistic $\mu/(\sigma/\sqrt n)$ of a period.

\begin{proposition}[The evidence of a measured halving]\label{prop:rs:half-life-decay-and-stability:power}
With prior probability $\pi$ that the IC truly halved, a measured halving (second-period mean at most half the first) has posterior probability
$\pi p_1/(\pi p_1 + (1 - \pi)p_0)$ of true decay, where $p_0$ and $p_1$ are the probabilities of the measurement without and with decay; for
normal means $p_1$ is close to $\frac12$ and $p_0$ falls with the period's $t$ statistic.
\end{proposition}

\begin{proof}
Bayes' rule on the two hypotheses. With decay, the second mean is centred on half the first's centre with the same noise, and the event
``second at most half the first'' is nearly symmetric about that centre, so $p_1 \approx \frac12$; without decay the event needs the noise to move
the means apart by half their level, whose probability falls as the level grows against the noise.
\end{proof}

With $\pi = \frac12$ and five-year periods (\cref{fig:rs:half-life-decay-and-stability:posterior}), a \omterm{def:rs:anatomy-of-a-predictor:predictor}{predictor} with the synthetic momentum's
noise (a five-year $t$ statistic of 0.72) that measures a halving has truly decayed with probability 0.57, barely more than the prior; one with
the momentum factor's pre-publication noise ($t = 1.93$), 0.72; one as steady as the synthetic surprise ($t = 7.6$), 0.999. A decline is
evidence in proportion to the \omterm{def:rs:anatomy-of-a-predictor:predictor}{predictor}'s precision; for noisy \omterm{def:rs:anatomy-of-a-predictor:predictor}{predictors}, the decision to cut must rest on other evidence (capacity, crowding,
the mechanism) or on longer samples.

\begin{omfigure}
\begin{tikzpicture}
\begin{axis}[width=10cm,height=5.4cm,xlabel={five-year $t$ statistic of the IC},ylabel={P(true decay $\mid$ measured halving)},
  tick label style={font=\small},label style={font=\small},xmin=0,xmax=8.3,ymin=0.45,ymax=1.02,grid=major,grid style={gray!25},
  table/col sep=comma]
\addplot[omDef,very thick,mark=*,mark size=1.4pt] table[x=t,y=posterior]{figdata/research/13-half-life-decay-and-stability/posterior.csv};
\node[font=\footnotesize,anchor=west] at (axis cs:0.9,0.53) {synthetic momentum};
\node[font=\footnotesize,anchor=west] at (axis cs:2.1,0.70) {momentum factor, 1963--1993};
\node[font=\footnotesize,anchor=east] at (axis cs:7.5,0.96) {synthetic surprise};
\end{axis}
\end{tikzpicture}

\omcaption{Posterior probability that a \omterm{def:rs:anatomy-of-a-predictor:predictor}{predictor} whose second five-year IC mean is at most half its first has truly halved, with a prior of one
half, against the five-year $t$ statistic of its IC (\cref{prop:rs:half-life-decay-and-stability:power}; 200\,000 simulated pairs per point).}\label{fig:rs:half-life-decay-and-stability:posterior}
\end{omfigure}

\section{Tutorial: dating a predictor's decline}\label{tut:rs:half-life-decay-and-stability:tut}

\begin{tutorial}
\textbf{Goal.} Measure the decay curves and half-lives of four synthetic \omterm{def:rs:anatomy-of-a-predictor:predictor}{predictors}, test a planted break with the CUSUM and sup-F tests, measure
three real factors before and after publication, and compute the evidence a measured halving carries. \textbf{End state:}
\cref{fig:rs:half-life-decay-and-stability:curves,fig:rs:half-life-decay-and-stability:break,fig:rs:half-life-decay-and-stability:factors,fig:rs:half-life-decay-and-stability:posterior};
the two tables.
\begin{tutsteps}
\item \textbf{Fit the curve} by least squares over a grid of half-lives, negative ICs included.
\omcode{code/firm/decay/firm_decay.py}{44}{56}{Least-squares half-life of an IC decay curve.}\label{lst:rs:half-life-decay-and-stability:fit}
\item \textbf{The CUSUM path and its boundary.}
\omcode{code/firm/decay/firm_decay.py}{85}{96}{Brown, Durbin and Evans's CUSUM for a constant mean.}\label{lst:rs:half-life-decay-and-stability:cusum}
\item \textbf{Run} \texttt{decay\_table()}, \texttt{stability()} with and without the planted break, \texttt{factors()}, \texttt{power()} and
\texttt{fig\_decay.py}; \texttt{rs\_fetch\_factors.py} once for the French statistics.
\end{tutsteps}
\textbf{What to change next.} Plant a gradual decay instead of a break (a drift that falls linearly over three years) and see which test finds it
first; repeat the post-publication table with the sample of each paper instead of 1963.
\end{tutorial}

\section{Build: the decay toolkit}\label{bld:rs:half-life-decay-and-stability:decay}

\begin{build}
\bfield{Purpose} The horizon, half-life and stability fields of every \omterm{def:rs:anatomy-of-a-predictor:card}{predictor card}, measured the same way, and monthly monitoring of the
\omterm{def:rs:anatomy-of-a-predictor:predictor}{predictors} in production.
\bfield{Interface} \texttt{ic\_decay(signal, ret, horizons, dates)}, \texttt{half\_life\_fit}, \texttt{rank\_autocorr}, \texttt{half\_life\_ar},
\texttt{turnover}, \texttt{block\_bootstrap}, \texttt{cusum}, \texttt{sup\_f}, \texttt{sup\_f\_critical}, \texttt{sup\_f\_pvalue},
\texttt{decline\_power(ic, sd, n, ratio)}; \texttt{firm.synthmkt} gains \texttt{pead\_break} and \texttt{pead\_after} (off by default).
\bfield{Rules} A decay curve uses single-period returns at each lag, never cumulative ones; every half-life is reported with the form it assumes;
a break test's p-value comes with the number of \omterm{def:rs:anatomy-of-a-predictor:predictor}{predictors} monitored.
\bfield{Acceptance tests} \texttt{code/firm/decay/tests/}: a planted five-day \omterm{def:rs:half-life-decay-and-stability:decay}{signal half-life} recovered by the curve fit and the autocorrelation; a
one-day signal; bootstrap shape; CUSUM silent on noise and crossing on a planted shift; sup-F locating the shift, its simulated critical value,
permutation p-values on a shift, on noise and on heavy-tailed noise; the power probabilities.
\bfield{Stretch} Bai--Perron for several breaks; a Bayesian change-point model with a prior on decay; decay curves by capacity bucket.
\end{build}

\begin{omsources}
\item R.~D.~McLean and J.~Pontiff, ``Does academic research destroy stock return predictability?'', \emph{Journal of Finance} 71(1), 2016.
\item T.~Chordia, A.~Subrahmanyam and Q.~Tong, ``Have capital market anomalies attenuated in the recent era of high liquidity and trading
activity?'', \emph{Journal of Accounting and Economics} 58(1), 2014.
\item K.~Hou, C.~Xue and L.~Zhang, ``Replicating anomalies'', \emph{Review of Financial Studies} 33(5), 2020.
\item R.~L.~Brown, J.~Durbin and J.~M.~Evans, ``Techniques for testing the constancy of regression relationships over time'', \emph{Journal of the
Royal Statistical Society B} 37(2), 1975.
\item D.~W.~K.~Andrews, ``Tests for parameter instability and structural change with unknown change point'', \emph{Econometrica} 61(4), 1993;
J.~Bai and P.~Perron, ``Estimating and testing linear models with multiple structural changes'', \emph{Econometrica} 66(1), 1998.
\item R.~W.~Banz, \emph{Journal of Financial Economics} 9(1), 1981; E.~F.~Fama and K.~R.~French, \emph{Journal of Finance} 47(2), 1992;
N.~Jegadeesh and S.~Titman, \emph{Journal of Finance} 48(1), 1993; Kenneth R. French Data Library.
\end{omsources}

\section{Exercises}

\begin{exercise}[$\star$]\label{exo:rs:half-life-decay-and-stability:1}
A signal's lag-one daily rank autocorrelation is 0.98. What are its half-life and the daily turnover of a book that holds it?
\end{exercise}

\begin{exercise}[$\star$]\label{exo:rs:half-life-decay-and-stability:2}
A factor earned 0.85\% a month before publication and 0.37\% after. What is the decline, and how does it compare with McLean and Pontiff's
average?
\end{exercise}

\begin{exercise}[$\star$]\label{exo:rs:half-life-decay-and-stability:3}
With $K = 94$ recursive residuals, where is the CUSUM's 5\% boundary at the first residual and at the last?
\end{exercise}

\begin{exercise}[$\star\star$]\label{exo:rs:half-life-decay-and-stability:4}
Every event drifts at a constant rate for 60 days after its announcement and the signal is observed at a uniformly random age
between 0 and 59 days. Show that the share of signals still drifting on day $t + h$ is $(60 - h)/60$, and explain why the \omterm{def:rs:half-life-decay-and-stability:decay}{IC decay curve} then
falls linearly to zero at 60 days rather than halving at a fixed rate.
\end{exercise}

\begin{exercise}[$\star\star$]\label{exo:rs:half-life-decay-and-stability:5}
A monthly IC has mean 0.04 and standard deviation 0.10. What is the five-year $t$ statistic, and roughly what posterior of true decay does a
measured halving give, with a prior of one half?
\end{exercise}

\begin{exercise}[$\star\star$]\label{exo:rs:half-life-decay-and-stability:6}
Why does a \omterm{def:rs:anatomy-of-a-predictor:predictor}{predictor} that loads on a common factor have a noisier monthly IC than one that does not, even with the same true information?
\end{exercise}

\begin{exercise}[$\star\star\star$]\label{exo:rs:half-life-decay-and-stability:7}
\emph{Coding.} With \texttt{rs\_decay.detection}, measure how often sup-F and the CUSUM detect a 50\% cut of the drift at year 7 over the
ten seeds, and place it between the chapter's removal and 70\% cut.
\end{exercise}

\begin{exercise}[$\star\star\star$]\label{exo:rs:half-life-decay-and-stability:8}
\emph{Find the flaw.} ``Our momentum signal's IC over the last 24 months is a third of its 10-year average and the CUSUM crossed its boundary last
month; we are switching it off.''
\end{exercise}

\section{Problem: Did It Die, or Was It Unlucky?}

\begin{problem}[{Weekend problem --- a decline, weighed}]\label{pb:rs:half-life-decay-and-stability:1}
The chapter's four synthetic \omterm{def:rs:anatomy-of-a-predictor:predictor}{predictors}, the surprise with its drift removed at year 7, and the French factors.

\medskip\noindent\textbf{Part I --- Horizons.}
\begin{enumerate}
\item Give the one-day IC and the \omterm{def:rs:half-life-decay-and-stability:decay}{signal autocorrelation} of each \omterm{def:rs:anatomy-of-a-predictor:predictor}{predictor}.
\item Why does the reversal's book turn over completely each day?
\item What half-life does an exponential fit give the surprise, and what is its planted life?
\item Why can momentum's and value's decay not be measured on single days?
\end{enumerate}

\medskip\noindent\textbf{Part II --- Years.}
\begin{enumerate}[resume]
\item Give the mean and standard deviation of the monthly IC for momentum and for the surprise.
\item What is the bootstrap interval for momentum's mean IC?
\item Why is momentum's IC so much noisier?
\item What do the CUSUM and sup-F tests say about momentum, which the simulator holds constant?
\end{enumerate}

\medskip\noindent\textbf{Part III --- Breaks.}
\begin{enumerate}[resume]
\item What are the surprise's mean IC before and after the planted removal?
\item Where does sup-F put the break, with what statistic and p-value? Why not exactly at month 48?
\item What does the CUSUM find?
\item What happens with a 70\% cut instead, in this market and over ten markets?
\item Give the three factors' declines after publication and the sup-F p-values over their whole histories.
\end{enumerate}

\medskip\noindent\textbf{Part IV --- The verdict.}
\begin{enumerate}[resume]
\item What are $p_0$ and $p_1$ for a \omterm{def:rs:anatomy-of-a-predictor:predictor}{predictor} with the synthetic momentum's noise?
\item State the \emph{named result}: the probability that a signal whose measured IC halved over five years has truly decayed, for the three
precisions of the chapter.
\item What $t$ statistic makes a measured halving 90\% convincing?
\item What other evidence would you seek before switching off a noisy \omterm{def:rs:anatomy-of-a-predictor:predictor}{predictor}?
\item How would you set a monitoring rule for fifty \omterm{def:rs:anatomy-of-a-predictor:predictor}{predictors} so that false alarms stay rare?
\item Which readings besides investors' trading can explain a post-publication decline?
\item In one sentence: when is a decline evidence?
\end{enumerate}
\end{problem}

\section{Interview questions}

\begin{interviewq}[$\star$ \iqroles{researcher}]\label{iq:rs:half-life-decay-and-stability:1}
What is the half-life of a signal, and how does it relate to turnover?
\end{interviewq}

\begin{interviewq}[$\star\star$ \iqroles{researcher}]\label{iq:rs:half-life-decay-and-stability:2}
How would you measure how fast a \omterm{def:rs:anatomy-of-a-predictor:predictor}{predictor}'s information decays with horizon?
\end{interviewq}

\begin{interviewq}[$\star\star$ \iqroles{researcher, trader}]\label{iq:rs:half-life-decay-and-stability:3}
Your best signal has had a bad year. How do you decide whether it is broken?
\end{interviewq}

\begin{interviewq}[$\star\star$ \iqroles{researcher}]\label{iq:rs:half-life-decay-and-stability:4}
Why do anomalies weaken after publication? Give three explanations and a way to tell them apart.
\end{interviewq}

\begin{interviewq}[$\star\star$ \iqroles{researcher}]\label{iq:rs:half-life-decay-and-stability:5}
Compare the CUSUM and sup-F tests for a change in a signal's mean.
\end{interviewq}

\begin{interviewq}[$\star\star\star$ \iqroles{researcher}]\label{iq:rs:half-life-decay-and-stability:6}
Derive the probability that a signal's second five-year mean IC is at most half its first when nothing has changed, as a function of the IC's
$t$ statistic.
\end{interviewq}
